\documentclass[11pt,a4paper]{article}
\usepackage[margin=2.2cm]{geometry}
\usepackage{amsmath,amssymb,booktabs,graphicx,xcolor,enumitem}
\usepackage[hidelinks]{hyperref}
\setlist{itemsep=2pt,topsep=3pt}
\newcommand{\ka}{\kappa_a}
\title{Coronary centerline tortuosity: choices and mathematics}
\author{Working document, ImageCAS-X development cohort}
\pdfstringdefDisableCommands{\def\ka{kappa_a}}
\date{26 September 2026 (provisional, not pre-registered)}

\begin{document}
\maketitle

\noindent\textbf{Status.} Every choice below was made on the ImageCAS-X training split (560 scans), which is a
development cohort: clinically referred patients, about half with plaque, mixed cardiac phase, no age or sex.
Nothing here is a normal range. No descriptor has yet been tested for agreement between reference and
CAS-Net-derived centerlines, and the region limits were chosen on the same scans they are evaluated on.
They must be frozen in a dated commit and checked on the validation split before any thesis claim.

\section{Input and preprocessing}

\begin{itemize}
\item \textbf{Data.} Only the delivered centerlines,
  \path{<id>.coronary_{left,right}_centerline.vtk}: LPS coordinates in mm, per-point
  \texttt{segment\_label}, native spacing about 0.45\,mm.
\item \textbf{Vessel path.} For a trunk label $t$ (RCA 9, LAD 2, LCx 3) the path is the longest geodesic
  path through edges whose two ends both carry label $t$. It is oriented from the ostium (RCA, via the
  \texttt{start\_points} array) or from the LM (LAD, LCx) to the distal end.
\item \textbf{No extra smoothing.} The dataset regenerated the centerlines by skeletonising the corrected
  masks and applying a Gaussian of $\sigma=0.5$\,mm over a 5-vertex window. The path is only resampled at
  $h=0.25$\,mm of arc length, giving points $\mathbf{x}_0,\dots,\mathbf{x}_n$ with arc length
  $s_i = \sum_{j<i}\lVert\mathbf{x}_{j+1}-\mathbf{x}_j\rVert$ and $L=s_n$. Adding $\sigma=1$\,mm changed
  $\ka$ medians by 2 to 4\,\% (Spearman 0.98 to 0.99). Any centerline from another source (CAS-Net,
  CGPS) must first pass through the same skeleton plus $\sigma=0.5$\,mm step.
\item \textbf{Exception.} Descriptors that need a second derivative (Section~\ref{sec:bishop}) use an extra
  Gaussian of $\sigma_0 = 1.25$\,mm along arc length, fixed from a noise bound before any cohort number.
\end{itemize}

\section{Descriptors}

\subsection{Arc/chord (distance metric)}
\[
  \mathrm{DM} = \frac{L}{\lVert \mathbf{x}_n - \mathbf{x}_0 \rVert}, \qquad
  \mathrm{DM}_{2D} = \frac{\sum_i \lVert P(\mathbf{x}_{i+1}-\mathbf{x}_i)\rVert}{\lVert P(\mathbf{x}_n-\mathbf{x}_0)\rVert},
\]
where $P\in\mathbb{R}^{2\times3}$ projects onto the image plane of an angiographic view. The viewing
direction for LAO/RAO angle $\alpha$ and cranial/caudal angle $\gamma$ is
$\mathbf{d} = (\sin\alpha\cos\gamma,\,-\cos\alpha\cos\gamma,\,\sin\gamma)$ in LPS, and $P$ has rows
$\mathbf{e}_1 = \mathbf{d}\times\hat{\mathbf{z}}/\lVert\cdot\rVert$ and $\mathbf{e}_2=\mathbf{d}\times\mathbf{e}_1$.
Standard views: RCA LAO\,30 ($\alpha=30^\circ$), LAD RAO\,20/CRA\,30, LCx RAO\,20/CAU\,30.
DM measures the overall course and grows with vessel length (Spearman 0.70 to 0.84 per whole vessel).

\subsection{Mean chord curvature $\ka$ (primary)}
Chord resampling places vertices $\mathbf{v}_0=\mathbf{x}_0,\mathbf{v}_1,\dots$ on the curve with
$\lVert\mathbf{v}_{k+1}-\mathbf{v}_k\rVert=\ell$ exactly (solved on each polyline segment). With unit
chords $\mathbf{u}_k = (\mathbf{v}_{k+1}-\mathbf{v}_k)/\ell$ and turning angles
$\theta_k = \arccos(\mathbf{u}_{k-1}\cdot\mathbf{u}_k)$ at vertex $\mathbf{v}_k$,
\[
  \ka^{(\ell)} = \frac{1}{L}\sum_k \theta_k \quad [\text{rad/mm}], \qquad \ell = 5\ \text{mm}.
\]
Total turning is $\Theta = \sum_k\theta_k = \ka L$. For a smooth curve $\theta_k \approx \int_{\text{chord}}\kappa\,ds$,
so $\ka$ is a low-pass mean absolute curvature at scale $\ell$. It is insensitive to vessel length (Spearman
0.2 to 0.3), not affected by dominance or image quality, and it cannot see a bend closer than one chord to
the path end (so it ignores an ostial hook by construction).

\subsection{Bends and the largest bend angle}
At $\ell=1$\,mm chords, a vertex is \emph{bending} if $\theta_k > \kappa_{\min}\ell$ with
$\kappa_{\min}=0.02$\,mm$^{-1}$ (bend radius under 50\,mm). A bend is a maximal run of bending vertices,
split where the binormal $\mathbf{b}_k=\mathbf{u}_{k-1}\times\mathbf{u}_k$ reverses
($\mathbf{b}_k\cdot\mathbf{b}_{k-1}<0$). Its angle is $\sum_{k\in\text{run}}\theta_k$ and its position the
$\theta$-weighted mean arc length. Reported: the \textbf{largest bend angle} per region. Bend counts
($\ge45^\circ$, $\ge90^\circ$) are descriptive only, because they depend on $\kappa_{\min}$ and smoothing.

\subsection{Bishop curvature spectrum}\label{sec:bishop}
On the $\sigma_0$-smoothed curve $\mathbf{c}(s)$ with unit tangent $\mathbf{T}$, carry a rotation-minimising
(Bishop) frame $(\mathbf{T},\mathbf{M}_1,\mathbf{M}_2)$ by parallel transport and write the curvature vector
$\mathbf{T}'$ as one complex number
\[
  \psi(s) = k_1 + i k_2 = \mathbf{T}'\cdot\mathbf{M}_1 + i\,\mathbf{T}'\cdot\mathbf{M}_2, \qquad |\psi| = \kappa .
\]
$\psi$ needs no torsion, is defined through inflections, and for a planar curve is the signed 2D curvature
up to a constant phase. This is what unifies the 2D and 3D descriptions.

\begin{itemize}
\item \textbf{Bending energy.} $\displaystyle k_B = \Big(\frac{1}{|R|}\int_R |\psi|^2\,ds\Big)^{1/2}$ (mm$^{-1}$),
  the RMS curvature over region $R$.
\item \textbf{Out-of-plane share.} With a Gaussian window $g_W$ of standard deviation $W=5$\,mm,
  \[
    f_{\text{twist}} = 1 - \frac{\int_R |g_W * \psi^2|\,ds}{\int_R g_W * |\psi|^2\,ds}\in[0,1].
  \]
  Squaring removes the sign, so $g_W*\psi^2$ is large only when bending keeps one plane over the window.
  $f_{\text{twist}}=0$ exactly for any planar curve, whatever its inflections; for a helix of curvature
  $\kappa$ and torsion $\tau$ it is $1-e^{-2\tau^2W^2}$.
\item \textbf{Wiggle share} (dropped, defined for completeness). An orthonormal DCT of $\psi$ is split by a
  power-complementary Gaussian, $|H_L(\omega)|^2 = e^{-(\omega s_c)^2}$, $|H_H|^2 = 1-|H_L|^2$,
  $s_c = \sqrt{\ln 2}\,\lambda_c/2\pi$, $\lambda_c=20$\,mm; $f_{\text{wiggle}} = \int_R|\psi_H|^2/\int_R(|\psi_L|^2+|\psi_H|^2)$.
\end{itemize}
The classical indices are moments of the same $\psi$ spectrum: RMS curvature is moment 0, a Grisan-type
turn index is about moment $-1$, and arc/chord is about moment $-2$, which explains why arc/chord follows
length and truncation.

\subsection{Regional attribution}
Every quantity is computed once on the whole vessel and then attributed to a region $R=[a,b)$ of arc length:
$\ka(R) = \sum_{k:\,s(\mathbf{v}_k)\in R}\theta_k/|R|$; bends are assigned by their centre; $k_B$,
$f_{\text{twist}}$ and $f_{\text{wiggle}}$ integrate their pointwise densities over $R$; DM uses the
sub-curve. Regional totals therefore add up to the whole-vessel value and no region loses its edges to a
chord or a filter. Regions shorter than 10\,mm, or without their landmark, give an explicit NaN row with the reason.

\section{Regions per vessel}
Take-off positions are the arc length of the trunk point joined by an edge to the first point of the side
branch label (nearest trunk point within 3\,mm if no edge exists).

\begin{center}
\small
\begin{tabular}{@{}lp{3.6cm}p{4.2cm}p{4.6cm}@{}}
\toprule
Vessel & Proximal & Mid & Distal / end window\\
\midrule
RCA & 5 to 40\,mm (first 5\,mm trimmed) & 40\,mm to $L-20$ & crux: last 20\,mm before R-PDA/R-PLA\\
LAD & LM to D1 & D1 to D2, else halfway to end & rest\\
LCx & LM to OM1 & OM1 to OM2, else halfway to end & rest\\
\bottomrule
\end{tabular}
\end{center}

\begin{itemize}
\item \textbf{RCA 5\,mm trim.} Many RCAs start with a sharp hook where the centerline leaves the aortic
  root. In scan 662 it holds most of the vessel's bending energy; trimming 5\,mm cut the proximal $k_B$
  from 0.121 to 0.041\,mm$^{-1}$. The LAD and LCx show no such hook (whole vs trimmed $\ka$, Spearman
  0.99 to 1.00), so they are not trimmed.
\item \textbf{Crux window.} Defined only when the RCA label ends at R-PDA/R-PLA: 512 of 513 right-dominant,
  6 of 15 co-dominant, 0 of 32 left-dominant scans.
\item \textbf{Landmark coverage.} Measurable proximal regions: RCA 560, LAD 424 (D1 median 18\,mm from the LM,
  127 regions under 10\,mm), LCx 365 (OM1 missing in 69 right-dominant scans, 126 under 10\,mm).
  Recommendation (untested for LAD and LCx): fixed proximal windows as primary, landmark segments as secondary.
\end{itemize}

\section{Evidence behind the choices (560 training scans)}

\textbf{Regional $\ka$ values are nearly independent}, so a single whole-vessel number hides where
bending sits (Spearman between regions):

\begin{center}
\small
\begin{tabular}{@{}lccc@{}}
\toprule
 & RCA (prox$_\text{trim}$/mid/crux) & LAD (prox/mid/distal) & LCx (prox/mid/distal)\\
\midrule
proximal vs mid & 0.55 & 0.13 & 0.14\\
proximal vs distal/crux & 0.21 & 0.20 & 0.10\\
mid vs distal/crux & 0.24 & 0.37 & 0.25\\
whole vs distal/crux & 0.43 & 0.89 & 0.71\\
\bottomrule
\end{tabular}
\end{center}
The whole-LAD value is mostly the distal (apical) LAD, which is not where first plaque forms.

\textbf{Added information.} Share of each descriptor's rank variance \emph{not} explained by
($\ka$, vessel length), in the plaque-relevant proximal and mid regions:

\begin{center}
\begin{tabular}{@{}lccc@{}}
\toprule
Descriptor & RCA & LAD & LCx\\
\midrule
$f_{\text{twist}}$ & 94 to 96\,\% & 98 to 99\,\% & 98 to 99\,\%\\
largest bend angle & 39 to 57\,\% & 69 to 83\,\% & 66 to 70\,\%\\
regional DM (3D) & 19 to 43\,\% & 64 to 66\,\% & 44 to 60\,\%\\
$k_B$ & 19 to 30\,\% & 48 to 69\,\% & 41\,\%\\
mean 10\,mm arc/chord & 16 to 17\,\% & 46 to 68\,\% & 42\,\%\\
\bottomrule
\end{tabular}
\end{center}
Mean 10\,mm arc/chord and $k_B$ agree with each other at Spearman 0.88 to 0.98 in every proximal and mid region, so at most
one of them can be kept. In short regions some of the apparent new information may be estimation noise.

\textbf{Confounds.} Proximal and mid $\ka$: no dominance effect (all $p\ge0.2$), $|\rho|\le0.09$ with image quality.
Whole-vessel DM shifts with dominance in all three vessels ($p$ from $10^{-5}$ to $10^{-10}$). Distal LCx
descriptors follow length ($\ka$ $\rho$ 0.47, bends 0.61), because the distal LCx extent depends on dominance.

\section{Decision}

\begin{center}
\small
\begin{tabular}{@{}lp{7.2cm}p{5.4cm}@{}}
\toprule
Status & Descriptor & Where / condition\\
\midrule
Primary & $\ka^{(5)}$ & every plaque-relevant region\\
Secondary & largest bend angle & same regions\\
Conditional & $f_{\text{twist}}$ & proximal and mid, only if reliable\\
Literature only & whole-vessel DM (3D and 2D view), whole-vessel $\ka$ & reported with length and dominance\\
Descriptive & bend counts $\ge45^\circ$; RCA crux; distal LCx & distal LCx stratified by dominance\\
Exploratory & $k_B$ & RCA only after the 5\,mm trim\\
Dropped & mean 10\,mm arc/chord, total turning, $f_{\text{wiggle}}$, regional 2D DM, regional bends $\ge90^\circ$,
  inflection count, ICM, full SOAM, writhe, box-counting dimension, SRVF distance to a line & \\
\bottomrule
\end{tabular}
\end{center}

\noindent\textbf{Open before adoption.}
\begin{enumerate}
\item Reliability per region: jitter and, above all, reference versus CAS-Net-derived centerlines.
  Whole-vessel $f_{\text{twist}}$ fell to ICC 0.15 to 0.32 under 0.5\,mm correlated jitter, while $\ka$
  held 0.74 to 0.81.
\item Fixed-window versus landmark proximal regions for the LAD and LCx.
\item Freeze regions and thresholds in a dated commit; confirm on the validation split.
\item CGPS needs a labeller (no artery names) and the same smoothing step before any of this applies.
\end{enumerate}

\section*{Figures}
\begin{center}
\includegraphics[width=\textwidth]{/work3/s254124/imagecasx_results/tortuosity_research/case_study/rca_tortuosity_comparison.png}\\[2pt]
{\small Two length-matched RCAs (scans 662, 518) in LAO\,30 with easy and advanced descriptors.}\\[8pt]
\includegraphics[width=0.85\textwidth]{/work3/s254124/imagecasx_results/tortuosity_research/case_study/regional_lao30.png}\\[2pt]
{\small The same RCAs split into proximal, mid and crux regions (before the 5\,mm trim is applied).}
\end{center}

\section*{Code and outputs}
{\small\raggedright
Scripts in \path{tortuosity/research/}: \path{simple/compute_simple.py}, \path{unified/bishop.py},
\path{case_study/easy.py}, \path{advanced.py}, \path{regional.py}, \path{regional_analysis.py}, \path{visualise.py}.
Outputs in \path{/work3/s254124/imagecasx_results/tortuosity_research/case_study/}:
\path{regional_train560.csv}, \path{regional_lad_train560.csv}, \path{regional_lcx_train560.csv} and the matching
\path{_analysis.txt} files.
Rebuild with \path{latexmk -pdf tortuosity_methods.tex}.\par}

\end{document}
